\section{Related Work and Novelty Boundary}

Low-rank optimization work spans at least three overlapping concerns: subspace
selection, discarded-gradient recovery, and optimizer-state compression. GaLore
and related projected optimizers reduce update rank directly; GoLore shows that
SVD-selected subspaces can be noise dominated in stochastic settings and uses
random low-rank projection to recover convergence guarantees. \Fira\ and
SubTrack++ add a full-rank recovery branch, but that branch is still driven by
projected adaptive information. Classical Error Feedback instead accumulates
compression error explicitly in a residual state.

LDAdam is the closest related method and the largest novelty risk. LDAdam
compensates low-dimensional Adam gradient and optimizer-state compression
errors through generalized Error Feedback and projection-aware state updates.
Our work addresses a different mechanism: the \Fira-style full-rank recovery
branch, where a discarded gradient component is reintroduced through norm-based
scaling. \BECR\ separates raw-gradient residual transmission from adaptive
recovery scaling in that branch. The \BECR-SGD theorem does not compensate Adam
optimizer-state compression, does not subsume LDAdam, and should not be read as
an Adam convergence theorem \citep{ldadam}.

Table~\ref{tab:related-work} in the appendix summarizes the boundary more
compactly. The short version is that this paper is about \emph{recovery-branch
mechanics}: when a norm-based full-rank correction is not Error Feedback, when
bounded recovery repairs a stale-subspace failure, and how far one can go with
an SGD abstraction before overlapping LDAdam's compressed-state theory.
